% Generated draft: requires mathematical and source review.
\documentclass[11pt]{article}
\usepackage[margin=25mm]{geometry}
\usepackage{fontspec}
\setmainfont{Latin Modern Roman}
\usepackage{amsmath,amssymb,mathtools,unicode-math}
\setmathfont{Latin Modern Math}
\usepackage{enumitem,needspace,xurl,hyperref}
\setlength{\parindent}{0pt}
\setlength{\parskip}{5pt}
\setlength{\emergencystretch}{3em}
\newcommand{\N}{\mathbb{N}}
\newcommand{\Z}{\mathbb{Z}}
\newcommand{\Q}{\mathbb{Q}}
\newcommand{\R}{\mathbb{R}}
\newcommand{\C}{\mathbb{C}}
\newcommand{\F}{\mathbb{F}}
\newcommand{\A}{\mathbb{A}}
\newcommand{\PP}{\mathbb{P}}
\newcommand{\E}{\mathbb{E}}
\newcommand{\eps}{\varepsilon}
\newcommand{\dd}{\,\mathrm{d}}
\begin{document}
\section*{UnsolvedMath ID: 30001055}
\textbf{Source problem code:} catalogue entry
\subsection{Definitions and mathematical statement}
A finite abstract simplicial complex on a vertex set $V$ is a family $\mathcal K\subseteq2^V$ closed under taking subsets. Its dimension is $\max_{\sigma\in\mathcal K}(|\sigma|-1)$; graphs give the one-dimensional case. A random simplicial complex is a random variable with values in such complexes, requiring a specified probability law. The source asks: Can one develop a theory of random simplicial complexes analogous to random graph theory? This is a broad high-dimensional combinatorics research program, not merely the existence of a probability distribution on finite complexes. No dimension regime, random model, independence rules, size parameter, properties to study, or criterion for an adequate analogy is supplied, so no particular threshold, topological limit, or classification theorem is substituted. The supplied 2026-08-21 triage labels the question open as a research program while explicitly recognizing substantial existing random-complex theory; it does not claim that no such theory or results exist.
\subsection{Short English statement}
Develop a theory of random simplicial complexes comparable to random graph theory. The source leaves the random models, dimensional regimes, and desired results open, and recognizes substantial existing work.
\subsection{Sources}
\begin{itemize}\item Local UnsolvedMath registry, permanent ID 30001055, source problem catalogue entry, \emph{Random Simplicial-Complex Theory}, proposed year 2008; no individual proposer supplied.\item Supplied citation: \emph{Discrete Geometry} (2009), Oberwolfach Reports, DOI 10.4171/owr/2008/44, pp.\ 2477--2556: \url{https://doi.org/10.4171/owr/2008/44}. Supplied literature label and attribution: Linial report.\item Literature assessment date: 2026-08-21. The status and progress claims are supplied dated curation, not independently verified here.\end{itemize}
\end{document}
